\chapter{代谢是一张网络，不是一条流水线}

\begin{kaichang}
晚饭时间到了。桌上一碗米饭、一盘炒鸡蛋、几筷子青菜、一块红烧肉。你先别动筷子——这一章的问题是：\textbf{三天之后，这顿饭在哪里？}

请在纸上写下你的猜测，最好带上比例。比如：三成变成了我的肌肉，四成变成能量烧掉了，两成变成脂肪存起来，剩下一成排出了体外。哪种去向最多？那块红烧肉，是不是更容易变成我身上的肉？

本章结束时我们回来对答案。你会发现，正确答案几乎一定是“每样都有一点”，而且去向不取决于这顿饭“是什么”，取决于你身体里一张巨大的、每时每刻都在自我调节的网。这张网，就是\textbf{代谢}。
\end{kaichang}

\section{现象层：这顿饭去了哪里}

先不谈分子，只看你能观察到的事实。

饭下肚半小时，饥饿感消失，身体微微发热——消化、吸收、转运这一整套动作本身就要花能量。饭后一小时，如果用血糖仪（许多糖尿病人家里都有）扎一下手指，会看到血糖读数升高；再过两小时，它又回到饭前水平。第二天早上称体重，数字和昨天差不太多：那顿饭的大部分已经悄悄离开了——混在你呼出的气里、汗液和尿液里，还有一部分压根没有被吸收，直接穿肠而过。

如果你连续几周吃得比消耗多，体重慢慢涨；反过来，慢慢降。剧烈运动时呼吸会加深加快；长时间不吃饭会头晕、心慌、注意力涣散。冬天呼出的“白气”，是代谢产生的水蒸气遇冷凝结的小水滴；尿液每天带走尿素和多余的盐分。这些全是看得见、测得到的现象。它们共同指向一个结论：\textbf{吃进去的东西既没有消失，也没有“原地变成你”，而是在不停地流动、转化、进出}。

想亲眼抓到“转化正在发生”的证据？不需要任何试剂，用你的呼吸就行。

\begin{shiyan}
\textbf{数一数身体工厂的“排烟速度”}

材料：一块有秒表功能的表或手机。

步骤：安静坐着休息三分钟后，数一分钟内自己的呼吸次数，记下来。然后原地快步走或做深蹲两分钟，停下后立刻再数一分钟的呼吸次数。休息五分钟，数第三次。

观察点：运动后呼吸次数增加了多少？几分钟后回到安静水平？呼出的这口气和安静时相比，多出来的东西是什么？

安全提示：量力而行，感到胸闷、头晕就立刻停下休息；饭后一小时内不要做。
\end{shiyan}

运动后呼吸加深加快，是因为你的肌肉正在加速“烧掉”燃料，而产生的废气——二氧化碳——必须从血液运到肺、更快地排出去。你刚才数的，其实是身体这座工厂的排烟速度。这座工厂从不熄火，区别只在转速。

\section{粒子层：两条方向相反的街道}

把镜头推到分子层面。那顿饭的主要“货物”是三类大分子：淀粉（第 46 章）、脂肪（第 47 章）和蛋白质（第 48 章）。它们个头太大，穿不过肠壁进入血液，必须先由酶剪成小零件：淀粉剪成葡萄糖，脂肪剪成脂肪酸和甘油，蛋白质剪成氨基酸。这些小零件被小肠吸收进入血液，随血流快递到全身每一个细胞。

进了细胞，小零件们走上两种方向的“街道”。

\textbf{分解代谢}：把分子拆小，沿途收获能量和电子。葡萄糖先在细胞质里被拆成丙酮酸（糖酵解，第 52 章），再进入线粒体被拆到只剩二氧化碳（柠檬酸循环，第 53 章）；拆下来的高能电子驱动 ATP 的合成。脂肪酸也被两个碳、两个碳地切段。方向是大变小，能量进账。

氨基酸拆下来还有一样麻烦的东西：含氮的氨基。氨对细胞有毒，不能直接扔进血液了事。肝脏把它改造成毒性低、易溶于水的尿素，交给肾脏随尿排出——这就是高蛋白大餐之后你会口渴、尿多的原因：身体在加班处理氮。你体检报告里的“尿素氮”一项，量的正是这条排污管的工作量。

\textbf{合成代谢}：花掉 ATP 提供的能量，把小分子拼成大分子。葡萄糖拼成糖原，存进肝脏和肌肉；脂肪酸重新拼成脂肪，存进脂肪组织；氨基酸拼成你自己的蛋白质——肌肉纤维、头发、酶、抗体。方向是小变大，能量支出。你长高的一厘米、伤口愈合长出的新皮、练了几个月才鼓起来的肱二头肌，全是合成代谢的施工成果。

请注意，这两条街不是互不往来的单行道。它们更像一张城市路网：同一片街区，有的车往东，有的车往西，路口彼此相连。本章剩下的内容，就是带你逛这些路口。

\begin{qiaoliang}
这张网的“车流量”有多大？用 ATP 算一笔账（ATP 的来历见第 50 章）。一个成年人正常活动一天，大约消耗 2000 千卡能量。在细胞内的真实条件下，水解 1 摩尔 ATP 大约能提供 12 千卡可用能量，所以身体每天要“花掉”约 $2000 \div 12 \approx 170$ 摩尔 ATP。1 摩尔 ATP 约重 507 克，折算下来，每天被合成又分解掉的 ATP 总质量约 $170 \times 507 \approx 8.6 \times 10^{4}$ 克——\textbf{超过一个成年人的体重}。然而你身体里任何时刻只存着大约 50 克（约 0.1 摩尔）ATP。结论只有一个：每个 ATP 分子每天平均要循环使用约 $170 \div 0.1 \approx 1700$ 次，不到一分钟就完成一次“放电、充电”。能量货币的意义不在于攒，而在于转。代谢首先是一场永不停歇的流动，其次才是一份库存清单。
\end{qiaoliang}

\section{规则层：所有大街在这里交汇}

第 52、53 章带你走过糖酵解和柠檬酸循环这两段主干道，现在站高一层俯瞰：糖、脂肪、蛋白质三条大街，其实在市中心汇入同一个环岛。

环岛的关键路口是一个叫\textbf{乙酰辅酶 A} 的小分子，可以把它想象成一辆装着两个碳原子的货运小车。葡萄糖拆到丙酮酸之后，丙酮酸变成乙酰辅酶 A；脂肪酸被逐段切碎，切下来的也是一截截乙酰辅酶 A；许多氨基酸脱掉含氮的氨基后，剩下的碳骨架也能变成乙酰辅酶 A，或者直接进入柠檬酸循环的中途站点。三种来源完全不同的货物，在这里混装。

小车们接下来有两个大方向：开进柠檬酸循环，彻底拆成 \chem{CO_2}，把电子交给呼吸链去造 ATP——这是“烧掉”；或者运往工地当原料——乙酰辅酶 A 是合成脂肪酸和胆固醇的起点，柠檬酸循环的中间产物是合成氨基酸的骨架，氨基酸再拼成蛋白质——这是“盖楼”。

路口是双向的：糖可以变成脂肪（所以米饭、面条吃多了同样长脂肪）；饥饿时，氨基酸可以转身变成葡萄糖去救急（这叫糖异生）。但每条路有自己的单行规则：脂肪酸的主体部分\textbf{不能}净转化为葡萄糖，所以长时间饥饿时，大脑会改用肝脏用脂肪造出的替代燃料——酮体。哪些路口能走、哪些不能，由每个反应的 $\Delta G$ 和有没有对应的酶决定（第 27、49 章），不是细胞“想”怎么走。

现在可以让符号登场了。葡萄糖被完全氧化的总反应式是：

\[\chem{C_6H_{12}O_6} + \chem{6O_2} \rightarrow \chem{6CO_2} + \chem{6H_2O}\]

这行式子看着眼熟——它和第 54 章光合作用的总式子像照镜子。但请回忆第 54 章的警告：两者走的路线完全不同，绝不互为“倒带”。也请把这行式子读成\textbf{总账}而不是\textbf{路径}：没有任何一个葡萄糖分子是被一步烧成灰的，它要走过几十步反应，每步都有专门的酶守着。总账告诉你起点、终点和能量收支，却不告诉你中间的路网。这正是本章标题想说的：黑板上这条式子像一条笔直的流水线，细胞里真实运行的是一张网。

\begin{center}
\begin{tikzpicture}
\node[draw, rounded corners] (glu) at (0,2.2) {葡萄糖、糖原};
\node[draw, rounded corners] (fa) at (0,0) {脂肪酸、甘油};
\node[draw, rounded corners] (aa) at (0,-2.2) {氨基酸};
\node[draw, rounded corners] (ac) at (4.4,0) {乙酰辅酶 A};
\node[draw, rounded corners] (tca) at (8.4,0) {柠檬酸循环};
\node (co2) at (11.8,0.9) {二氧化碳、ATP};
\node (fat) at (8.4,-2.2) {脂肪、胆固醇};
\node (pro) at (11.8,-0.9) {新的蛋白质};
\draw[->] (glu) -- (ac);
\draw[->] (fa) -- (ac);
\draw[->] (aa) -- (ac);
\draw[->] (ac) -- (tca);
\draw[->] (tca) -- (co2);
\draw[->] (ac) -- (fat);
\draw[->] (aa) -- (pro);
\draw[<->] (glu) -- (aa);
\end{tikzpicture}
\end{center}

上图是人为的表示法：每个方框代表的不是一个分子，而是时刻在进出的一大群分子；箭头表示“\textbf{能}转化”，不表示“此刻正在转化”。方框的大小、位置也和真实细胞毫无关系。

\section{规则层：没有总指挥，网络怎么不乱}

一座城市只有街道、没有红绿灯，早就瘫痪了。你的每个细胞都是一座代谢小城，全身几十万亿座，却没有任何总指挥。它靠什么不乱？

第一套装置是\textbf{反馈抑制}。看一条典型的合成路线：原料 A 经过 B、C，最终变成产物 D。当 D 库存充足，D 会亲自“回头”，结合在这条路线\textbf{第一步}的那个酶上，让它停工。这像面包店：货架上的面包堆满了，老板直接通知面粉厂停止送货，而不是等面包烂在店里。开关设在第一步而不是最后一步，是因为这样最省——中间产物 B、C 压根不会白白堆积起来。

第二套装置藏在 D 结合的方式里。D 结合的位置不在酶的活性位点上，而在酶表面的另一个“口袋”——一卡进去，酶的整体形状改变，活性位点跟着变形，底物就进不去了。这种“在远处结合、改变形状、调节活性”的手法叫\textbf{变构调节}（蛋白质的形状会变这一点，第 48 章已经埋过伏笔）。分子不用堵在门口，打个电话就能让机器改设置。

第三套装置：能量状态本身就是信号。ATP 充足时，ATP 会直接变构抑制糖酵解的关键酶——“钱够了，葡萄糖先别拆了”；ADP、AMP 积累时，同一个酶被激活——“快没钱了，加速”。同一个分子既是货币又是仪表盘读数。这类调节不靠激素、不靠大脑，每个细胞内部每秒钟都在发生——第 52 章说过“代谢途径受能量状态调节”，现在你看到了它的分子手法。

用符号把反馈抑制写成一句话：设路线为 $A \rightarrow B \rightarrow C \rightarrow D$，则 D 会抑制催化 $A \rightarrow B$ 的那一步。箭头向右，信号向左——一进一出，路线自己就稳在了合适的产量上，不需要谁下命令。

\section{证据层：科学家怎样知道身体在“不停翻新”}

\begin{zhengju}
二十世纪三十年代之前，主流的图景是“房子加炉灶”：身体是已经盖好的房子，食物是燃料，烧掉供能；组织里的蛋白质被认为基本稳定，磨损了才修补一点。

打破这幅图景的工具，你在第 54 章刚见过——同位素示踪（同位素的概念见第 7 章）。三十年代初，科学家分离出氢的重同位素氘（比普通氢多一个中子）：化学行为几乎不变，却能在仪器里被认出来。哥伦比亚大学的舍恩海默（Rudolf Schoenheimer）和同事里滕伯格（David Rittenberg）立刻意识到它的用途：给食物分子贴上这种看不见的“标签”，追踪它们在动物体内的去向。

第一批实验：把含氘的脂肪酸喂给小鼠。几天后检测，标记没有全部烧掉，而是大量出现在小鼠储存的脂肪里——哪怕小鼠的体重没变。这意味着储存脂肪不是“进了仓库就不动的囤货”，而是在不停地被存入、取出、更新。第二批实验更彻底：用带标记的氮-15 标记某种氨基酸喂给动物，几天后，全身各处的蛋白质——肝脏、肌肉、皮肤——都带上了标记；甚至标记还出现在了\textbf{别的}氨基酸分子里，说明氨基酸之间也在互相转换。

会不会是假象？比如，标记只是吸附在组织表面？研究者把动物组织研磨、用化学方法把蛋白质真正分离出来再测——标记嵌在分子内部，洗不掉。会不会只有肝脏这样“活跃”的器官才更新？肌肉、皮肤的蛋白质同样带着标记，只是速度慢一些。

结论在 1942 年被写成一本影响深远的书《身体成分的动态状态》：\textbf{身体的建筑材料处于不停拆解与重建的动态之中}，体重不变只是“拆”和“建”恰好平衡，而不是什么都没发生。后来更精细的测量又修正了细节：不同组织的翻新速度天差地别——肝脏的蛋白质几天就换掉一半，肌肉的蛋白质以周计，而眼球晶状体里的蛋白质几乎终生不换。动态是普遍的，速度各不相同。今天医院里查幽门螺杆菌的“碳-13 呼气试验”，用的正是同一套贴标签的思想。
\end{zhengju}

\section{规则层：全身的红绿灯}

单个细胞能自我调节，但你的肝脏、肌肉、脂肪组织、大脑是分工不同的器官，谁告诉它们“现在该存还是该取”？这套全身协调靠激素——用血液快递的化学消息。代谢领域最重要的两个信使是\textbf{胰岛素}和\textbf{胰高血糖素}，都由胰腺分泌，作用恰好相反。它们怎样把消息送进细胞内部，第 58 章会细讲；这一章只看它们指挥什么。

\textbf{刚吃完饭}：血糖升高，胰岛素出动。它通知肌肉和脂肪细胞打开葡萄糖通道，把血糖收进来；通知肝脏和肌肉把葡萄糖拼成糖原库存；通知脂肪细胞把多余的原料合成脂肪。血糖随之回落。一句话：丰收了，入仓。

\textbf{饭后几小时到夜间}：血糖回落，胰高血糖素出动。它通知肝脏把糖原拆回葡萄糖、释放入血，保证大脑和红细胞的供应（它们几乎只认葡萄糖）。一句话：开仓放粮——但只动肝脏这个粮仓，肌肉里的糖原是肌肉的私产，不外卖。

\textbf{饥饿超过一天}：肝糖原（只有约 100 克，不到一天的口粮）见底，身体转向脂肪。脂肪组织放出脂肪酸供全身燃烧；肝脏把一部分脂肪酸改造成酮体，给大脑当替代燃料。同时有少量肌肉蛋白质被拆解，拆出的氨基酸经糖异生补充葡萄糖——这就是长期饥饿时肌肉和脂肪一起减少的原因。

\textbf{运动}：肌肉先烧自己的糖原；几分钟后，血液运来的葡萄糖和脂肪酸加入供能；运动时间越长，脂肪供能的比例越高。注意，是“比例”在变——脂肪从第一分钟就参与供能，并不存在“运动半小时后才开始燃烧脂肪”的开关。你的肌肉更像混合动力车，电和油时刻都在混用，只是配比随路况调整。

这套红绿灯一旦失灵，就是病。1 型糖尿病患者的胰腺造不出胰岛素：丰收的消息发不出去，葡萄糖堵在血液里进不了细胞——血糖高得危险，细胞却在挨饿，患者需要终身注射胰岛素。另一个例子更能说明“网络”的价值：苯丙酮尿症（PKU）患者体内，处理氨基酸苯丙氨酸的那个酶坏了，一个路口堵死，上游堆积的苯丙氨酸会损伤发育中的大脑。对策不是吃药修酶，而是饮食管理——少吃含苯丙氨酸的食物，让车流不超过这个路口的通行能力。今天新生儿的足跟血筛查常规包含 PKU：一张网络图，直接变成了救命的饮食清单。这也再次提醒：营养和医学问题讲究剂量、证据和个体差异，任何“某种食物好或坏”的断言，都要先问“对谁、多少量”。

\section{边界层：这张地图什么时候会失灵}

网络模型很强大，但用过头就会骗人。记住四个边界。

第一，本章画的箭头只表示“能转化”，不表示“此刻正在转化”。同一张图，在吃饱的你和饥饿的你体内，通量分布完全不同——图是静态的，网是活的。

第二，平面图把一切都画在了一起，真实的细胞却把不同通路关在不同的“房间”里：糖酵解在细胞质，柠檬酸循环在线粒体。同一个分子身处哪个房间，命运就不同。空间安排是网络的重要一层，第 57 章专门讲。

第三，图上每个“路口”（酶）在不同人身上版本不同。乳糖不耐受（第 46 章）、苯丙酮尿症，都是某个路口的通行能力天生不同。所以不存在对所有人同样“正确”的饮食——同一张网，运行参数因人而异。

第四，网络知识目前还做不到精确预言“你吃某样东西一定会怎样”。验血能看到的血糖、胆固醇等指标，只是整张网的几个窗口，医生据此做判断时看的也是趋势和整体，而不是单个数字。任何向你保证“能精确算出你该吃什么”的说法，都超出了今天的证据。

\begin{wuqu}
“吃核桃补脑，吃猪皮美容，吃腰子补肾——吃什么补什么。”

这是把代谢当成直通流水线的典型后果：仿佛食物从嘴到目标器官之间有一条专线。真实的旅程是：猪皮的胶原蛋白先在消化道里被拆成氨基酸，汇入全身的氨基酸“大池子”；接下来它们可能被拼成你的头发、你的酶、你伤口处的新皮肤，也可能被拆掉氨基、烧成二氧化碳，或者转身变成葡萄糖和脂肪。没有任何一个氨基酸分子“记得”自己来自猪皮，也没有任何机制把它定向发往你的脸。同理，核桃里的脂肪大部分被燃烧或重新合成储存，并不会排队进入大脑。

反例就在眼前：严格素食者从不吃一块肌肉，照样长出肌肉——原料是豆子里的氨基酸，图纸在身体里。这类说法里唯一站得住的，只剩“原料清单”：少数必需氨基酸和必需脂肪酸身体自己造不出来，必须从食物获得（第 46、47 章）。食物提供的是砖，身体决定盖什么——施工图从来不在食物上。
\end{wuqu}

\section{再看那碗饭}

现在可以对答案了。三天之后，那碗米饭的碳原子们大约这样分布：一部分已经变成二氧化碳，从你的肺里呼了出去；一部分还存着，是肝脏和肌肉里的糖原，或者脂肪细胞里新合成的脂肪；还有一小部分兜兜转转，成了某个蛋白质里的碳骨架。那块红烧肉的蛋白质拆成了氨基酸，汇入大池子，可能去修补你的肌肉，也可能被烧掉——它并没有“专线直达”你的肌肉。

最关键的是：这些比例不固定。同一碗饭，遇上饭后睡觉的你和饭后打球的你，命运分布截然不同。决定去向的是网络当时的状态——库存多少、激素发什么消息、你在做什么。没有哪一粒分子“注定”去哪，只有概率和调节。你当初猜的比例，只要不是“全去一个地方”，就都猜对了一半。

\section{继续深挖：这张网住在哪里}

走到这里，这条因果链已经很长：原子、分子、反应、有机分子、生物大分子，再到本篇讲的——分子怎样被组织成一张自我供能、自我维修、自我调节的代谢网络。第六篇到此结束。

但你可能已经憋着一个问题：我们一直说“细胞里”“线粒体里”，这张网到底运行在一个什么样的容器里？几十步反应、几千种分子挤在一个比头发丝还细的空间里，为什么不乱成一锅粥？细胞不是烧杯——它有内部结构、有运输系统、有通讯机制，更像一座城，而不是一个袋子。下一篇我们就进城：第 56 章先问一个最大的难题——这样的化学城市，最初是怎样从一堆分子中自发搭起来的？第 57 章再带你看城市内部的分区与物流。你会发现，本章这张平面的网，将被折进一座立体的、会呼吸的城。

\begin{tiaozhan}
稳态不等于平衡——这是理解代谢最容易混淆的一对概念。第 30 章讲的化学平衡，是正逆反应速率相等、没有\textbf{净}流动的状态；而血糖的稳态，是流入等于流出、有持续流动的状态。模型是一只同时开着水龙头和排水口的浴缸：水位不变，但水在不停更换。代入数字感受一下：你的血液约 5 升，血糖浓度约每升 1 克，总共约 5 克葡萄糖——这就是“水位”。空腹时肝脏以每小时几克到十几克的速度向血液放糖，同时各组织以同样的速度消耗。若流量是每小时 6 克，那么池中一个葡萄糖分子平均只待 $5 \div 6 \times 60 = 50$ 分钟就被换掉。水位几乎不动，水却早已换了几轮——这就是“动态状态”的定量含义。再进一步：若流入突然超过流出（比如一顿大餐），水位上升本身会触发调节（胰岛素）把流出开大——反馈让浴缸有了自动水位控制。能用“流入、流出、库存”三个量描述的系统，都可以用这套直觉分析：体温、血钙，甚至一个城市的货物存量。跳过本框不影响主线。
\end{tiaozhan}

\section*{练一练}

\begin{sikaoti}
\item 画一张物质流图：一块红烧肉里的脂肪和蛋白质，从被吃下到抵达“市中心环岛”，各经过哪些拆解步骤？从环岛出发，至少画出三条不同的去向，并标注每条是“烧掉”“入仓”还是“盖楼”。
\item 用反馈抑制的思路解释：为什么“大量吃某种氨基酸，身体就会大量合成对应的蛋白质”不成立？多出来的氨基酸最终去了哪里？
\item 把血糖画成一只浴缸的水位：标出“进水口”和“出水口”各是什么，并说明一顿大餐之后，胰岛素怎样扳动这些开关，让水位不至于涨得太高。
\item 一个人连续三天只喝水、不吃饭。按时间顺序写出他身体燃料来源的变化，并解释：为什么脂肪和肌肉会同时减少？
\item 本章算过：每个 ATP 分子每天循环约 1700 次。请算出它平均多少秒完成一次循环（一天约 86000 秒），并说说这个数字说明“能量货币”和纸币的一个本质不同。
\item 你的亲戚说：“我最近皮肤不好，要多吃猪蹄补一补。”用本章的模型写三句话回应他——既要指出错在哪里，也要承认哪一点点道理（如果找得到的话）。
\end{sikaoti}
